\begin{frame}{Änderungsanträge (DROI)}
\tiny
\vspace{-0.45cm}
\renewcommand{\arraystretch}{0.9}
\setlength{\tabcolsep}{2pt}
\setlength{\abovetopsep}{0pt}
\setlength{\belowrulesep}{0pt}
\begin{tabularx}{\textwidth}{@{}>{\centering\arraybackslash}m{1.55cm} >{\RaggedRight\arraybackslash}m{\dimexpr\linewidth-1.55cm-2\tabcolsep\relax}@{}}
\toprule
\multicolumn{1}{m{1.55cm}}{\centering\textbf{Fraktion}} & \multicolumn{1}{m{\dimexpr\linewidth-1.55cm-2\tabcolsep\relax}}{\centering\textbf{Änderungsantrag}} \\
\midrule
\includegraphics[width=2cm,height=1cm,keepaspectratio]{Bilder/EVP.png} & ii) Liegt die Anerkennungsquote des Herkunftslandes unter \textcolor{red}{\sout{20}}\textcolor{green}{\uline{12}} \% der geflüchteten Personen oder besteht ein relevantes Sicherheitsrisiko, greift das \\
\includegraphics[width=2cm,height=1cm,keepaspectratio]{Bilder/SD.png} & \textcolor{red}{\sout{iii) In begründeten Einzelfällen kann das Asylschnellverfahren auch in Drittstaaten stattfinden.}} \\
\includegraphics[width=2cm,height=1cm,keepaspectratio]{Bilder/RE.png} & iii) \textcolor{red}{\sout{In begründeten Einzelfällen kann das}}\textcolor{green}{\uline{Das}} Asylschnellverfahren \textcolor{green}{\uline{kann }}auch in Drittstaaten stattfinden. \\
\includegraphics[width=2cm,height=1cm,keepaspectratio]{Bilder/PfE.png} & ii) Liegt die Anerkennungsquote des Herkunftslandes unter \textcolor{red}{\sout{20}}\textcolor{green}{\uline{40}} \% der geflüchteten Personen oder besteht ein relevantes Sicherheitsrisiko, greift das \\
\includegraphics[width=2cm,height=1cm,keepaspectratio]{Bilder/Green.png} & i) Die Dauer des Asylschnellverfahrens beträgt maximal \textcolor{red}{\sout{12}}\textcolor{green}{\uline{20}} Wochen und umfasst eine individülle Prüfung der gestellten Anträge. Jeder antragstellenden Person \\
\bottomrule
\end{tabularx}
\end{frame}

